\section*{Abstract}

AI-enabled workflows fail in ways that are frequently mischaracterized --- as model quality problems when they are organizational problems, as one-off hallucinations when they are structural verification gaps, and as isolated output defects when they are propagating failures with no single point of origin. This report develops a working framework for reasoning about where error enters an AI-enabled workflow, how it propagates, how it can be detected and contained, and what a defensible standard of success actually looks like. It proceeds from a purpose statement (comparative reliability against a strong human baseline, not perfection) through job definition, system decomposition, specification and verification difficulty, single-output diagnosis, and into hazard analysis, propagation topology, assurance argument, and operational resilience. The framework is a lens, not a certification procedure: it is meant to make a specific deployment decision more defensible, not to replace the judgment of the people accountable for it. A closing section records open questions the framework does not resolve.
